\begin{abstract}
Trimming methods for robust clusterwise regression discard a fixed fraction of the data. Too low
a level breaks the fit; too generous a level can trim away a small group. We first propose
a group-recovery step that can follow any trimming or flagging method: it searches the discarded
units for a line, tests whether those near it form a peak rather than a band, and restores the
line as a group when the likelihood of Gaussian groups plus uniform noise improves by a margin
like that of the Bayesian information criterion. In simulations it repaired the failures of a
generous TCLUST-REG level with unequal groups, but not with three or four groups. The second
proposal, ESF (exact-subsample flagging), is a flagging procedure without a trimming level: it solves the clusterwise least-squares
problem exactly on small subsamples, flags units far from the best fit, and draws later
subsamples from the rest. Two constants stand in for the level: a subsample size, set from a
lower bound on the smallest group proportion, and a cap on the flagged set. It is meant for
data about whose contamination nothing is known: TCLUST-REG at the fixed level 0.30, followed by
reweighting and the recovery step, was as accurate as ESF on average up to a fifth of outliers, and
higher levels were more accurate beyond. The
flagged fraction estimates the contamination only when the errors are close to Gaussian. On taxi
fares with known tariffs, ESF found both in every sample.
\end{abstract}
